\documentclass[french,11pt]{book}
\usepackage{teach}
\usepackage[french]{babel}
\usepackage{amsmath,amssymb}
\usepackage{array}
\newcommand{\R}{\mathbb{R}}
\begin{document}
\begin{center}
{\Large \textbf{Devoir surveillé -- Série statistiques à deux variables}}\\[2mm]
Terminale technologique
\end{center}

\section*{Devoir surveillé}
\begin{enumerate}
  \item Le tableau suivant donne le rang $x$ d'une année et un chiffre d'affaires $y$ en milliers d'euros : $(1;42)$, $(2;45)$, $(3;49)$, $(4;53)$, $(5;56)$, $(6;61)$.
  \item Représenter le nuage de points et déterminer une équation de droite d'ajustement affine à la calculatrice.
  \item Estimer le chiffre d'affaires pour le rang 9.
  \item Expliquer pourquoi cette estimation doit être utilisée avec prudence.
\end{enumerate}

\end{document}
